%\section{Mindmap}
%\vfill
\begin{figure}[!t]
	\centering
	\begin{tikzpicture}[
		mindmap,
		text width=2.5cm,
		align=center,
		every node/.style={concept},
		concept color=gray!30,
		level 1/.append style={
			level distance=5cm,
			minimum size=3cm,
			text width=2.8cm,
			font=\large,
		},
		level 2/.append style={
			level distance=3.5cm,
			minimum size=2cm,
			text width=1.8cm
		},
		]
		
		\node[minimum size=4.5cm, text width=4.3cm, line width=6pt, fill=gray!10] (TTT) { \Large\textbf{PVPS} }
		child[concept color=blue!10,grow=210] { 
			node {\textbf{Revisão bibliográfica}} 
			child[grow=75] { node {Métodos de estimação} }
			child[grow=120] { node {Estimação com restrições} }
			child[grow=240] { node {Tempo de trânsito} }
		}
		child[concept color=red!10, grow=270] {
			node {\textbf{\ingles{Travel-Time Tomography}}}
			child[grow=270] { node {Aplicações} }
			child[grow=15] { node {Adaptações} }
			child[grow=225] { node {Possíveis expansões (SVD)} }
			child[grow=330] {node {Outros temas} }
		}
		child[concept color=green!10, grow=0] {
			node {\textbf{Análise estatística}} 
			child[grow=105] {
                node {Bancos de dados}
                % child[grow=]
            }
			child[grow=60] { node {\ingles{Framework} avaliativo} }
            child[grow=-60] { node{\ingles{Packages} de manipulação de dados} }
			% child[grow=15] { node {Modelo inicial \\(3)} }
			% child[grow=-30] { node {Outros \\(3)}}
		}
		;	
	\end{tikzpicture}
	\caption{\textsl{Mindmap} da \cref{sec_rssfi}.}
	\label{fig:mindmap_sec4}
\end{figure}